\chapter{Cryptographic assumptions}

The security theorems of CatCrypt are \emph{reductions}: each takes the hardness
of a standard primitive as a hypothesis and concludes protocol security. This
chapter collects those hypotheses. Every assumption is phrased as an
\emph{advantage} --- a distance between a real and an ideal game, or a winning
probability in a search game --- and the assumption itself states that this
quantity is bounded (negligible) for all adversaries. All games are expressed in
the probabilistic program monad, and each advantage is measured with the
distinguishing notion of Chapter~3.

\section{Bilinear pairing groups}

\begin{definition}[Bilinear pairing group]
  \label{def:pairing-group}
  \lean{CatCrypt.Crypto.PairingGroup}
  \leanok
  A \emph{pairing group} (Type~III) bundles three cyclic groups
  $G_1, G_2, G_T$ of common prime order $p$, generators $g_1 \in G_1$,
  $g_2 \in G_2$, and a pairing $e : G_1 \to G_2 \to G_T$ that is bilinear in
  each argument ($e(a\cdot b, y) = e(a,y)\cdot e(b,y)$ and
  $e(x, a\cdot b) = e(x,a)\cdot e(x,b)$) and non-degenerate ($e(g_1, g_2) \neq 1$).
  Each group is finite with $|G_1| = |G_2| = p$, and $g_1$, $g_2$ generate their
  groups.
\end{definition}

\begin{lemma}[Pairing bilinearity on exponents]
  \label{def:pairing-bilinear}
  \uses{def:pairing-group}
  \lean{CatCrypt.Crypto.PairingGroup.e_bilinear}
  \leanok
  For integer exponents $a, b$,
  $e(g_1^{\,a},\, g_2^{\,b}) = e(g_1, g_2)^{\,a b}$. This is the identity that
  makes pairing-based verification equations checkable from public group
  elements.
\end{lemma}

\begin{definition}[KZG structured reference string]
  \label{def:kzg-srs}
  \uses{def:pairing-group}
  \lean{CatCrypt.Crypto.PairingGroup.srs₁}
  \leanok
  For a secret $\alpha \in \mathbb{Z}_p$ and degree bound $t$, the KZG
  structured reference string in $G_1$ is the tuple
  $\mathrm{srs}_1(\alpha, t) = (g_1, g_1^{\alpha}, g_1^{\alpha^2}, \dots,
  g_1^{\alpha^t})$; the companion $\mathrm{srs}_2(\alpha) = (g_2, g_2^{\alpha})$
  in $G_2$ supplies the verification key.
\end{definition}

\section{Discrete-logarithm assumptions}

\begin{definition}[Discrete logarithm advantage]
  \label{def:dl}
  \uses{def:pairing-group, def:advantage}
  \lean{CatCrypt.Crypto.Assumptions.DL_Advantage}
  \leanok
  In the DL game the challenger samples $x \in \mathbb{Z}_p$, hands the
  adversary $h = g_1^{\,x}$, and the adversary wins if it returns $x$. The
  \emph{DL advantage} of $A$ is $\Pr[\text{$A$ outputs }x]$. The DL assumption
  states this is negligible: no efficient adversary recovers exponents.
\end{definition}

\begin{definition}[Diffie--Hellman group]
  \label{def:ddh-def}
  \lean{CatCrypt.Crypto.Assumptions.DDHDef}
  \leanok
  Abstractly, a Diffie--Hellman group is a scalar-multiplication operation
  $\mathrm{ecdh} : W \to W \to W$ together with a (heap-independent) key
  generator producing pairs $(\text{private}, \text{public})$. This models
  $g \mapsto g^a$ and key sampling without committing to a concrete group.
\end{definition}

\begin{definition}[Decisional Diffie--Hellman advantage]
  \label{def:ddh}
  \uses{def:ddh-def, def:advantage}
  \lean{CatCrypt.Crypto.Assumptions.DDH_Advantage}
  \leanok
  The DDH advantage is the distinguishing distance between the real game, which
  gives the adversary $(g^a, g^b, g^{ab})$, and the ideal game, which gives
  $(g^a, g^b, c)$ for a fresh uniform $c$. The DDH assumption states this
  advantage is negligible: the true shared secret is indistinguishable from a
  random group element.
\end{definition}

\begin{definition}[Computational Diffie--Hellman advantage]
  \label{def:cdh}
  \uses{def:ddh-def}
  \lean{CatCrypt.Crypto.Assumptions.CDH_Advantage}
  \leanok
  In the CDH game the adversary receives $(g^a, g^b)$ and wins by outputting the
  shared secret $g^{ab}$. The CDH advantage is $\Pr[\text{$A$ outputs }g^{ab}]$;
  the CDH assumption states it is negligible. CDH is the computational analogue
  of DDH: distinguishing (DDH) implies computing (CDH), but not conversely.
\end{definition}

\begin{definition}[co-CDH advantage]
  \label{def:cocdh}
  \uses{def:pairing-group}
  \lean{CatCrypt.Crypto.Assumptions.CoCDH_Advantage}
  \leanok
  In a Type~III pairing group the co-CDH game samples $a, b \in \mathbb{Z}_p$,
  gives the adversary $(g_1^{\,a}, g_2^{\,b})$ across the two source groups, and
  the adversary wins by outputting $g_1^{\,ab}$. The co-CDH advantage is this
  winning probability. It is the hardness assumption underlying BLS signatures.
\end{definition}

\begin{definition}[Gap Diffie--Hellman advantage]
  \label{def:gapdh}
  \uses{def:ddh-def, def:advantage}
  \lean{CatCrypt.Crypto.Assumptions.GapDH_Advantage}
  \leanok
  GapDH strengthens CDH by granting the adversary a DDH decision oracle: it must
  still \emph{compute} $g^{ab}$ from $(g^a, g^b)$, even though it can \emph{check}
  candidate triples. The advantage is the distinguishing distance between the DH
  output and a fresh random value; the assumption states it is negligible. GapDH
  is the standard assumption for X25519-based protocols (WireGuard, Noise), and
  any DDH group is a fortiori a GapDH group.
\end{definition}

\begin{definition}[$t$-Strong Diffie--Hellman advantage]
  \label{def:tsdh}
  \uses{def:kzg-srs}
  \lean{CatCrypt.Crypto.Assumptions.tSDH_Advantage}
  \leanok
  In the $t$-SDH game the challenger samples $\alpha$ and publishes the SRS
  $(g_1, g_1^{\alpha}, \dots, g_1^{\alpha^t})$. The adversary wins by producing a
  pair $(c, h)$ with $h = g_1^{\,1/(\alpha + c)}$, checked via the pairing
  equation $e(h, g_2^{\,\alpha + c}) = e(g_1, g_2)$. The $t$-SDH advantage is
  this winning probability; it is the core hardness assumption for KZG
  polynomial commitments.
\end{definition}

\begin{definition}[$q$-SDH group]
  \label{def:qsdh-group}
  \lean{CatCrypt.Crypto.Assumptions.qSDH_Group}
  \leanok
  An abstract cyclic group of prime order carrying scalar arithmetic
  (addition, multiplication, inverse) and scalar multiplication $g \mapsto g^a$,
  finite and with decidable equality. This is the algebraic setting for $q$-SDH,
  abstracting the $G_1$ component of a pairing group without the pairing itself.
\end{definition}

\begin{definition}[$q$-Strong Diffie--Hellman advantage]
  \label{def:qsdh}
  \uses{def:qsdh-group}
  \lean{CatCrypt.Crypto.Assumptions.qSDH_Advantage}
  \leanok
  The challenger samples a secret $x$ and hands the adversary the challenge
  $(g, g^x, g^{x^2}, \dots, g^{x^q})$; the adversary wins with a pair $(c, h)$
  satisfying $h^{\,x + c} = g$, i.e.\ $h = g^{\,1/(x+c)}$. The $q$-SDH advantage
  is this winning probability. It underlies BBS/BBS+ signatures and anonymous
  credentials, capturing the same hardness as $t$-SDH in a pairing-free setting.
\end{definition}

\section{Oracle Diffie--Hellman}

\begin{definition}[Oracle Diffie--Hellman advantage]
  \label{def:odh}
  \uses{def:ddh-def, def:advantage}
  \lean{CatCrypt.Crypto.Assumptions.ODH_Advantage}
  \leanok
  ODH augments a DH group with a keyed hash $H$ whose key is the DH secret. The
  advantage is the distinguishing distance between the real game, which returns
  $H(g^{ab}, \ell)$ for a sampled label $\ell$, and the ideal game, which returns
  a fresh random value. The ODH assumption states this is negligible; taking $H$
  to be the identity recovers DDH.
\end{definition}

\begin{lemma}[DDH and PRF imply ODH]
  \label{def:odh-reduction}
  \uses{def:odh, def:ddh, def:prf}
  \lean{CatCrypt.Crypto.Assumptions.odh_of_ddh_prf}
  \leanok
  If the DH secret can be replaced by a random key at cost $\varepsilon_{ddh}$
  (a DDH step) and the hash keyed by a random value is a PRF at cost
  $\varepsilon_{prf}$, then the ODH advantage is bounded by
  $\varepsilon_{ddh} + \varepsilon_{prf}$, via a two-hop hybrid through the
  triangle inequality.
\end{lemma}

\section{Hash-function assumptions}

\begin{definition}[Collision-resistance advantage]
  \label{def:cr-search}
  \uses{def:advantage}
  \lean{CatCrypt.Crypto.Assumptions.CR_Search_Advantage}
  \leanok
  For a hash $h$, the search-CR game has the adversary output a pair
  $(x_1, x_2)$ and wins when $h(x_1) = h(x_2)$ with $x_1 \neq x_2$. The CR
  advantage is this winning probability; collision resistance states it is
  negligible.
\end{definition}

\begin{definition}[Target collision resistance advantage]
  \label{def:tcr}
  \uses{def:advantage}
  \lean{CatCrypt.Crypto.Assumptions.TCR_Advantage}
  \leanok
  In the TCR (second-preimage) game the adversary receives a random target $x$
  together with $h(x)$ and must find $x' \neq x$ with $h(x') = h(x)$. The TCR
  advantage is this winning probability. TCR is weaker than CR: collision
  resistance implies target collision resistance.
\end{definition}

\begin{definition}[Preimage-resistance advantage]
  \label{def:pre}
  \uses{def:advantage}
  \lean{CatCrypt.Crypto.Assumptions.Pre_Advantage}
  \leanok
  In the preimage game the adversary receives $y = h(x)$ for a random $x$ and
  wins by outputting any $x'$ with $h(x') = y$. The Pre advantage is this
  probability. For compressing hashes it is implied by collision resistance.
\end{definition}

\section{One-way functions and generators}

\begin{definition}[One-wayness advantage]
  \label{def:owf}
  \uses{def:advantage}
  \lean{CatCrypt.Crypto.Assumptions.OWF_Advantage}
  \leanok
  For a function $f$, the inversion game samples $x$, hands the adversary
  $f(x)$, and it wins by returning any preimage $x'$ with $f(x') = f(x)$. The
  OWF advantage is this winning probability; one-wayness states it is negligible.
\end{definition}

\begin{definition}[One-way permutation]
  \label{def:owp}
  \uses{def:owf}
  \lean{CatCrypt.Crypto.Assumptions.OWPDef}
  \leanok
  A one-way permutation is a one-way function that is additionally a bijection.
  RSA is the canonical conjectured instance.
\end{definition}

\begin{definition}[Pseudorandom-generator advantage]
  \label{def:prg}
  \uses{def:advantage}
  \lean{CatCrypt.Crypto.Assumptions.PRG_Advantage}
  \leanok
  For a stretch function $G$, the PRG advantage is the distinguishing distance
  between $G(\text{seed})$ for a random seed and a uniformly random output. The
  PRG assumption states this is negligible; a secure PRG is in particular a
  one-way function.
\end{definition}

\begin{definition}[RSA inversion advantage]
  \label{def:rsa}
  \uses{def:advantage}
  \lean{CatCrypt.Crypto.Assumptions.RSA_Advantage}
  \leanok
  The RSA game samples a key pair and a random preimage $x$, sets
  $y = \mathrm{rsa}(pk, x)$, and asks the adversary for $x'$ with
  $\mathrm{rsa}(pk, x') = y$. The RSA advantage is the distinguishing distance
  between this real game and the always-reject ideal game; the RSA assumption
  states it is negligible.
\end{definition}

\section{Symmetric primitives}

\begin{definition}[Pseudorandom-permutation advantage]
  \label{def:prp}
  \uses{def:advantage}
  \lean{CatCrypt.Crypto.Assumptions.PRP_Advantage}
  \leanok
  For a block cipher $E$, the PRP advantage is the distinguishing distance
  between the keyed permutation $E(k, \cdot)$ (random key) and a random
  \emph{function}. Together with the switching bound below this yields the full
  PRP security bound.
\end{definition}

\begin{definition}[Strong-PRP advantage]
  \label{def:sprp}
  \uses{def:advantage}
  \lean{CatCrypt.Crypto.Assumptions.SPRP_Advantage}
  \leanok
  The strong-PRP game grants the adversary both the forward oracle $E(k, \cdot)$
  and the inverse oracle $D(k, \cdot)$; the SPRP advantage is the distinguishing
  distance from a pair of independent random functions. This is the notion
  needed by modes that decrypt (e.g.\ CBC).
\end{definition}

\begin{definition}[Pseudorandom-function advantage]
  \label{def:prf}
  \uses{def:advantage}
  \lean{CatCrypt.Crypto.Assumptions.PRF_Advantage}
  \leanok
  For a keyed function $F$, the (single-query) PRF advantage is the
  distinguishing distance between $F(k, x)$ for a random key and a uniformly
  random output. The PRF assumption states this is negligible.
\end{definition}

\begin{definition}[PRP/PRF switching bound]
  \label{def:prp-prf-switch}
  \lean{CatCrypt.Crypto.Assumptions.PRP_PRF_SwitchingBound}
  \leanok
  The birthday term $q(q-1)/(2|W|)$ bounding the statistical gap between a random
  permutation and a random function over $q$ queries on a domain of size $|W|$.
  It is the additive correction relating the PRP advantage (which uses a random
  function as ideal) to true permutation security.
\end{definition}

\begin{definition}[MAC forgery advantage]
  \label{def:mac}
  \uses{def:advantage}
  \lean{CatCrypt.Crypto.Assumptions.MAC_Advantage}
  \leanok
  In the (indistinguishability form of) MAC game a random key is sampled and the
  adversary receives the tag on its chosen message; the MAC advantage is the
  distinguishing distance between this real tag and a uniformly random tag. The
  unforgeability assumption states it is negligible --- equivalent to
  single-query strong unforgeability up to a $1/|W|$ term.
\end{definition}

\begin{definition}[AEAD confidentiality advantage]
  \label{def:aead}
  \uses{def:advantage}
  \lean{CatCrypt.Crypto.Assumptions.AEAD_Advantage}
  \leanok
  For an authenticated-encryption scheme the (IND-CPA form) AEAD advantage is the
  distinguishing distance between the real game, which returns
  $\mathrm{encrypt}(k, \mathit{pt}, \mathit{ad})$ under a random key, and the
  ideal game, which returns a uniformly random ciphertext. Security states this
  is negligible.
\end{definition}

\begin{definition}[AEAD IND-CCA2 advantage]
  \label{def:aead-cca}
  \uses{def:aead}
  \lean{CatCrypt.Crypto.Assumptions.AEAD_CCA_Advantage}
  \leanok
  The IND-CCA2 game equips the adversary with both an encryption oracle and a
  decryption oracle; the advantage is the distinguishing distance between real
  encryption-plus-decryption and (random encryption, reject-all decryption). This
  combines confidentiality and ciphertext integrity, and implies the IND-CPA
  notion above.
\end{definition}
